\documentclass[UTF8,fontset=none,11pt,a4paper]{ctexart}
\usepackage[left=21mm,right=21mm,top=27mm,bottom=22mm,headheight=26pt,headsep=6mm]{geometry}
\usepackage{fontspec}
\setmainfont{texgyretermes-regular.otf}[BoldFont=texgyretermes-bold.otf,ItalicFont=texgyretermes-italic.otf]
\setCJKmainfont{Songti.ttc}[Path=/System/Library/Fonts/Supplemental/,FontIndex=6,AutoFakeBold=2]
\setCJKsansfont{Songti.ttc}[Path=/System/Library/Fonts/Supplemental/,FontIndex=6,AutoFakeBold=2]
\usepackage[table]{xcolor}
\definecolor{TableHead}{HTML}{EEF0ED}
\definecolor{Muted}{HTML}{666C68}
\usepackage{booktabs,longtable,array,fancyhdr,needspace,xurl}
\usepackage[unicode,hidelinks]{hyperref}
\urlstyle{same}
\linespread{1.32}
\setlength{\parindent}{2em}
\setlength{\parskip}{3pt}
\setlength{\LTpre}{8pt}\setlength{\LTpost}{10pt}
\ctexset{section={format=\fontsize{15}{20}\selectfont\bfseries,beforeskip=18pt,afterskip=8pt},subsection={format=\fontsize{12}{16}\selectfont\bfseries,beforeskip=12pt,afterskip=6pt},subsubsection={format=\normalsize\bfseries,beforeskip=9pt,afterskip=4pt}}
\pagestyle{fancy}\fancyhf{}
\renewcommand{\headrulewidth}{.3pt}
\renewcommand{\footrulewidth}{0pt}
\fancyhead[L]{\parbox[b]{\linewidth}{\fontsize{8}{10}\selectfont\color{Muted}退休账户转移：五年持有与两层费用·退休账户顾问项目披露资料}}
\fancyfoot[C]{\small\color{Muted}\thepage}
\hypersetup{pdftitle={退休账户顾问项目披露资料},pdfsubject={仅供模拟使用·退休账户转移：五年持有与两层费用·03-01},pdfauthor={金融决策模拟材料}}
\begin{document}
{\noindent\fontsize{20}{25}\selectfont\bfseries 退休账户顾问项目披露资料\par}
\vspace{7pt}{\noindent\small\color{Muted}仅供模拟使用·退休账户转移：五年持有与两层费用·03-01\par}\vspace{10pt}
\section*{项目服务与管理权限}
机构M通过投资顾问项目提供账户组合管理。项目使用不同风险程度的模型组合。根据客户问卷中提供的风险承受能力、投资期限、策略偏好及其他信息，客户取得列有资产配置和具体投资的提案，此后按适当组合管理资产。市场情况、客户需要及配置调整等因素可导致模型组合变化。


组合投资。项目从托管平台可用的基金中选择关联和非关联基金。平台可用范围受到购买和出售附加收费等条件限制。机构排除产生额外交易成本的基金或类别可降低自己的成本，但也可能排除表现、费率或其他指标更有利的基金。项目通过披露处理该利益冲突。


\textbf{份额选择。}同一共同基金不同份额投资同一底层证券组合，具有相同投资目标或政策，但费用、支出和初始最低投资额可能不同。项目通常采用其可用范围内面向机构的类别；基金未提供、项目不符合资格或未获最低额豁免时可采用其他类别。


机构类别通常不收12b-1费，但仍可能支付其他分销、行政、子过户或持有人服务费。机构从基金收到的部分费用按项目政策记作账户抵扣；其他服务商自行收到并保留的费用不由机构返还。项目会定期检查使用机构类别的资格，并在可操作时自行决定转换，不保证客户始终持有基金中最有利的类别。


\textbf{再平衡。}项目每天监控组合相对目标配置及权重的偏离。市场变动、存款或提款使配置出现达到项目标准的显著偏离时，符合余额要求的账户会再平衡。项目对超配使用比例阈值，对低配使用金额阈值，并考虑现金积累情况。标准和方法由机构或次级顾问决定，可自行调整。


\textbf{投资限制。}客户可口头或书面要求不投资某基金或用替代基金。但项目不接受不符合策略、指引、投资理念或运作方式的限制，也不能控制基金本身持有的底层证券。是否接受限制由机构决定。限制可令账户表现不同于推荐模型，整体回报可能较低；获接受限制列入提案和随后季度绩效报告。客户可以联系顾问修改或取消。


转入后保留的已有资产。客户可申请保留已有证券，须在开户问卷或转入前书面指定，并满足投资标准和合资格保留清单。持仓可能受到个别证券、行业或资产类别集中度限额。已有资产计入项目计费，其存在可能影响组合表现和再平衡。以后不再合资格时，可能不事先通知即出售或退回，出售可产生税务、递延销售费或其他成本。


\textbf{现金归集。}账户现金按开户申请所列归集选项处理。归集安排与可用银行可变化，存款利率可能低于项目以外其他现金安排。现金计入确定费用档位的账户价值，但项目费用不针对现金余额扣收。


\section*{项目成本}
项目费。参加项目支付按资产计算的项目费。适用费表依开户日期及已有项目客户或家庭关系不同。本费表适用于满足披露条件的新账户；其他账户应查阅自己的顾问协议。一般条款允许提前30日书面通知后变更费用，通知后继续参加视为同意。


\Needspace{5\baselineskip}
\begingroup\small\setlength{\tabcolsep}{4pt}\renewcommand{\arraystretch}{1.35}
\begin{longtable}{>{\centering\arraybackslash}m{0.69167\dimexpr\linewidth-4\tabcolsep\relax}>{\centering\arraybackslash}m{0.30833\dimexpr\linewidth-4\tabcolsep\relax}}
\toprule
\rowcolor{TableHead} \centering\arraybackslash \textbf{项目资产价值中\mbox{的部分}} & \centering\arraybackslash \textbf{\mbox{年费率}} \\
\midrule\endfirsthead
\toprule
\rowcolor{TableHead} \centering\arraybackslash \textbf{项目资产价值中\mbox{的部分}} & \centering\arraybackslash \textbf{\mbox{年费率}} \\
\midrule\endhead
\midrule\endfoot
\bottomrule\endlastfoot
\centering\arraybackslash 最初150,000美元 & \centering\arraybackslash 1.15\% \\ \addlinespace[3pt]
\centering\arraybackslash 接下来的150,000美元 & \centering\arraybackslash 1.00\% \\ \addlinespace[3pt]
\centering\arraybackslash 接下来的450,000美元 & \centering\arraybackslash 0.85\% \\ \addlinespace[3pt]
\centering\arraybackslash 接下来的250,000美元 & \centering\arraybackslash 0.75\% \\ \addlinespace[3pt]
\centering\arraybackslash 接下来的500,000美元 & \centering\arraybackslash 0.70\% \\ \addlinespace[3pt]
\centering\arraybackslash 接下来的1,500,000美元 & \centering\arraybackslash 0.65\% \\ \addlinespace[3pt]
\centering\arraybackslash 接下来的1,000,000美元 & \centering\arraybackslash 0.60\% \\ \addlinespace[3pt]
\centering\arraybackslash 接下来的1,000,000美元 & \centering\arraybackslash 0.50\% \\ \addlinespace[3pt]
\centering\arraybackslash 超过5,000,000美元\mbox{的部分} & \centering\arraybackslash 0.40\% \\
\end{longtable}\endgroup
混合费表对账户价值的各段适用不同费率。账户进入较高档位时，仅该档位内资产使用较低费率，由此产生适用于账户的加权费率。


\textbf{现金余额。}不论采用哪份费表，现金归集余额均计入判断费档的合计价值，但计算项目费时排除账户现金余额。


\textbf{家庭合并。}希望采用家庭合并时应通知顾问。该方法将家庭成员项目账户价值相加，以达到较高规模档位并降低加权费率。本人、配偶、父母、子女或同住人员持有或受益的合资格项目账户可以合并。合并后每个家庭取得一份共同季度绩效报告，并授权与指定家庭内其他账户持有人共享本人账户绩效资料；该授权不允许他们在本人账户交易。欲停止合并，应通知顾问。


\textbf{两层费用。}项目费不含底层基金固有的投资顾问、行政、分销、过户、托管、法律、审计及其他费用，也不包含基金投资相关的递延销售费或赎回费。因此客户直接承担项目费，并作为基金持有人间接承担基金费用。某些IRA或ERISA合资格账户投资关联基金时，可依条款取得关联管理费抵扣。费用总额可能高于或低于机构其他投资项目。


\textbf{12b-1和其他基金费用。}机构从基金收到的相应费用按项目政策抵扣至客户账户；其他服务商独立取得并保留的款项不由机构返还。基金收费具体见说明书及补充资料。


\textbf{项目费所含服务。}费用覆盖顾问或顾问团队、账户管理、建议开发、资产托管、通过项目经纪渠道交易、客户报告以及项目共同基金交易产生的赎回费。客户另外要求的服务和其他经纪账户专项交易费用按相应安排处理。


费用扣收。项目费按适用年率及每日投资资产价值累计，季度结束后扣取。扣费可使用现金或出售基金份额。开户或终止不足一个完整季度，按规则比例计算。

\end{document}